\section{The Rule in Operation}
\label{sec:applications}

No single fact pattern displays both front-door functions equally well. \emph{Garcia} is the deployment-record application: the operator is visible, but the deployment record is not. \emph{Mobley v. Workday, Inc.} is the divided-stack application: identity, possession, legal control, and technical control occupy different entities. The cited orders are used only for the procedural posture and record-allocation problems they document, not to predict any later merits result.

\subsection{An Integrated Deployment: \emph{Garcia}}
\label{sec:garcia-application}

\emph{Garcia} begins where the rule is easiest to serve. Character Technologies operated the visible consumer application. The complaint identified the company, its founders, and Google entities and described months of interaction through a known service. The difficult questions concerned the record behind the interface: which versions and policies governed, how features were allocated among the application and its components, what relational and self-harm risks were evaluated, what warnings and interventions appeared, and which actors could change access or operation.

The published decision supplies a careful procedural anchor. Taking well-pleaded allegations as true for Rule 12 purposes, the district court treated Character.AI as plausibly alleged to be a product and allowed central product and negligence theories to proceed. It declined on that record to treat the generated sequence as protected speech where defendants had not identified the relevant expressive choice or speaker. The order made no finding of defect, causation, or liability, and the later settlement supplied no merits determination.\lrfootnote{See \casecite[1165--69, 1174--84]{garcia2025}; \casecite[1]{garciaDismissal2026}; \casecite[1]{garciaLien2026}.}

\subsubsection{Trigger and First Answer}

The alleged event and pleaded facts illustrate a notice that, if verified, would satisfy the threshold for a relational deployment that a sectoral designation could cover. The complaint identified legally cognizable injuries, including wrongful death, and linked them to sustained interaction with the service. It challenged concrete features: access by a minor, anthropomorphic presentation, continuity and memory, relational engagement, warnings, and asserted omissions in safety intervention. The requested records would bear on recognized negligence, products, warning, and consumer theories rather than a freestanding claim to inspect a chatbot.

Character Technologies would receive the front-door notice as visible operator. Stage One would identify the application and model versions active for Sewell's account; whether the requested covered functions ran; warnings, detections, blocks, escalations, and interventions recorded for the event; the entities that supplied or configured relevant components and the custodians and legal controllers of the listed records; and the roles with update, restriction, rollback, suspension, or termination authority.

The complaint's allegations concerning the founders and Google demonstrate why identification should remain functional. It alleged personal direction by the founders and technical, financial, and organizational contributions by Google.\lrfootnote{See First Amended Complaint \P\P~24--25, 50--72, \emph{Garcia v. Character Techs., Inc.}, No. 6:24-cv-01903-ACC-EJK (M.D. Fla. Nov. 9, 2024), ECF No. 11; \casecite[1165--66, 1183--84]{garcia2025}.} Position, investment, former employment, or infrastructure supply alone would not establish material control. The first answer would identify whether any enterprise retained a risk-specific right over the model, application configuration, evaluation, release condition, update, or stopping decision. That record tests the control lane without presuming it from corporate association.

\subsubsection{Relational Features Define the Risk-Specific Record}

The production request can be narrower because relational research identifies provider-governed variables. Studies of companion discontinuity and user reports connect continuity, relationship framing, dependence, and recurring interaction patterns to experienced harms.\lrfootnote{See \bluecite[11--25]{defreitasReplika2025}; \artcite[3547--49, 3555--69]{banksDeletion2024}; \artcite[5926--38]{laestadiusTooHuman2024}; \bluecite[5--19, 22--23]{zhangDarkSide2025}.} A preregistered four-week study likewise found increased separation distress and dependency-risk profiles under specified relationship-seeking and repeated-exposure conditions.\lrfootnote{See \bluecite[2--3, 5--8, 10--11]{kirkSteering2026}. The preprint does not establish commercial-population incidence, historical foreseeability, clinical injury, or tort causation in \emph{Garcia}.}

The legal contribution lies in the variables rather than a universal effect estimate. Relationship seeking, continuity, memory, dose, age, crisis signal, and provider intervention are features an operator can evaluate and govern. They specify the record relevant to intended use, notice, alternative design, warning, monitoring, and causal sequence. Table~\ref{tab:relational-record} maps each feature to the risk-specific record it requires.

\begin{table}[htbp]
\centering
\caption{Relational Features as Control-Record Factors}
\label{tab:relational-record}
\small
\begin{tabularx}{\textwidth}{>{\raggedright\arraybackslash}p{0.28\textwidth} X}
\toprule
Deployment feature & Risk-specific record \\
\midrule
Persistent persona, memory, proactive messages, affection, or exclusivity & Intended relational use; exposure and continuity testing; memory policy; age-specific evaluation; monitored dependence signals. \\
Claims of companionship, therapy, expertise, autonomy, or capacity to act & Target audience; product requirements; substantiation; warning design; expected reliance; escalation responsibility. \\
Engagement objectives tied to relational behavior & Metric ownership; optimization target; risk constraint; release decision; evaluation of lower-exposure alternatives. \\
Self-harm or crisis classification & Detection version; threshold; validated population; false-positive and false-negative analysis; response policy; incident-specific trigger and intervention. \\
Access by minors & Age rule and enforcement; minor-specific evaluation; parental or guardian features where applicable; restriction and termination authority. \\
Discontinuity, account intervention, or service shutdown & Continuity notice; export or transition function; crisis safeguard; rollback and account-restriction decision. \\
\bottomrule
\end{tabularx}
\end{table}

The table also disciplines marketing evidence. Calling a service a companion or therapist does not make the model a legal person and does not establish defect. It can identify the use the enterprise cultivated, the exposure it expected, and the evaluations responsive to that design. A generic ``not human'' notice enters the same record with its timing, presentation, audience, and relation to the service's features.

\subsubsection{Stage Two and the Merits}

After Stage One established the versions, roles, and incident sequence, Stage Two could reach age-access and relational-risk evaluations; self-harm policies and configuration history; pre-event complaints and incident patterns; release and change approvals; escalation and response; and available account restrictions, interruptions, or rollbacks. Time, version, population, feature, and risk would bound the request. Unrelated conversations and user data remain excluded or minimized, and source code or weights require a showing that documentation, event logs, and representative testing cannot resolve a material issue.

The same record could favor the defense. The papers disputed character selection, role play, user-edited messages, platform rules, content blocking, and an in-chat fiction notice.\lrfootnote{See \casecite[1167--69]{garcia2025}.} Versioned configuration and event fields could show which features came from a third-party creator, which controls fired, which warnings appeared, and whether an alternative would have changed the sequence. Governing law would then address classification, foreseeability, feasibility, and causation; later research or safety measures could identify questions without proving historical knowledge, admissibility, or earlier feasibility.\lrfootnote{See Fed. R. Evid. 407; \emph{Garcia}, 785 F. Supp. 3d at 1169.}

\subsection{A Divided Deployment: The Vendor Hosts the Data; the Employer Controls It}
\label{sec:mobley-application}

\emph{Mobley} presents the complementary architecture. Applicants encounter employer portals while a platform vendor supplies processing and recommendation features. Employers can own or legally control applicant data and decide how a feature affects hiring; the vendor can host the data and control products, versions, defaults, and some evaluation. The cited orders expose four recurring record-allocation problems that the front door must solve without predicting discrimination: feature identity, divided legal control, privilege, and safe exit.

\subsubsection{Feature Identity and the Event Path}

The first problem is not whether ``AI'' was generally available. It is which function processed a particular application and what reached the employer. At the pleading stage, the district court treated allegations that employer customers delegated traditional hiring functions to Workday as sufficient to proceed on a statutory agency theory. The same order required a different theory to connect a particular employer's conduct, Workday's knowledge, and the assistance it allegedly supplied.\lrfootnote{See \casecite[806--08, 815--16]{mobley2024}. The order granted in part and denied in part a motion to dismiss; it made no finding of discriminatory operation or ultimate agency.} Nominal identity was visible. The operative chain among employer, feature, output, use, and knowledge was not.

Later notice-stage proceedings illustrate the fields. Candidate Skills Match, when enabled by a customer, parsed a job description and applicant materials, extracted skills, and assigned a categorical match. Workday described customers as able to enable, disable, use, or ignore features. The court reserved what outputs meant and how customers used them for later adjudication while permitting targeted discovery to identify potential collective members.\lrfootnote{See \casecite[2, 4, 10--13, 18--20]{mobleyCollective2025}. Preliminary certification authorized notice and preserved Workday's right to seek decertification; it was not a merits or final class ruling.}

For a designated employment-screening deployment, the employer receiving the application is the visible operator. Stage One asks the applicant for the position, application date, adverse disposition, covered portal, and the best available event key---a requisition number, application ID, tenant URL, or timestamp. It does not require the applicant to name the vendor feature or allege that a hidden score actually participated. The first answer resolves that question.

The employer's event receipt identifies products active for the application; versions and enabled features; fields, criteria, and preferences the employer supplied; any score, sort, rank, screen, threshold, or recommendation received; the downstream action; and whether no automated function processed the event. The custodian map identifies the supplier, account roles, record locations, legal controllers, and contractual retrieval, update, and suspension rights. A linked supplier answer resolves the same identifier within its lane: tenant and requisition, product and version, whether a function ran, whether it generated and transmitted an output, which settings came from the customer, and which entities held the listed records and authority roles.

Together, those answers separate five propositions: a feature was available; the customer enabled it; it processed the event; it generated or transmitted an output; and the customer used that output. None follows from the preceding proposition, and establishing the sequence does not establish that the output caused an unlawful decision.

\subsubsection{Physical Hosting Versus Legal Control}

The second problem is custody. A 2026 discovery order described customer applicant data as stored on Workday servers while the agreement placed ownership of customer content with the customer and allowed the customer to seek an injunction against disclosure. Applicants had subpoenaed thirty-five customers, and some customers pointed back to Workday as the efficient source. On that record, the court held that physical storage did not establish Workday's legal right to obtain and produce the customer data on demand under Rule 34.\lrfootnote{See \casecite[1--3, 8--10]{mobleyDiscovery2026}. The holding concerns the agreement and record before the court; it does not classify every hosted Workday record.}

The front door does not reverse that holding. It changes deployment design before the dispute. The employer keeps the control sheet, position definition, customer settings, application, downstream review, disposition, notice, and event identifier. The supplier keeps version lineage, product documentation, defaults, run and transmission logs, and the operational evaluation record assigned to its lane. The predeployment agreement gives the employer a right to retrieve defined supplier fields and authorizes a coordinated hold for customer-controlled content. Both receive the same deployment and event identifiers, preventing a platform label and customer label from preserving incompatible slices.

That architecture distinguishes technical hosting from statutory responsibility. A covered supplier subject to the Act answers its assigned lane. Another supplier responds through the employer's retrieval right. The applicant no longer has to choose between the host and legal controller before seeing the contract and custodian map; Stage One identifies both, and ordinary Rule 34 or Rule 45 analysis governs any later production beyond the statutory minimum.

\subsubsection{Privilege Versus the Operational Minimum Record}
\label{sec:mobley-privilege}

The third problem is the line between legal advice and operation. The same discovery order held particular Candidate Skills Match bias-testing materials privileged because counsel selected the data and used the tests to provide legal advice. It separately ordered production of Workday's EEO-1 and federal-contractor compliance materials based on their relevance to knowledge of potential demographic disparities.\lrfootnote{See \casecite[3, 7--8, 10--11]{mobleyDiscovery2026}. The order protected identified attorney-directed testing; it did not classify every evaluation or underlying business fact as privileged.}

The front door preserves that distinction. Attorney communications and qualifying counsel-directed analyses remain protected. The operational minimum---feature identity, version, intended use, validation owner, deployed threshold, event output, complaints, monitoring, and change history---exists because the covered system operated. A privilege log can identify protected analysis while underlying nonprivileged data and ordinary-course records follow existing law.\lrfootnote{See \casecite[395--96]{upjohn1981}; \statcite{frcp26WorkProduct}.} The statute therefore protects privilege without allowing an enterprise to satisfy its deployment-record duty only inside a legal-advice file.

\subsubsection{Safe Exit and Coordinated Stage Two}

The fourth problem is ending the special proceeding for the wrong lane. A 2025 order concerning HiredScore directed Workday to identify customers that enabled specified features and supplied a concrete exclusion rule: a customer could leave the preliminary collective if Workday could determine that it received no score or rank or did not use the feature to screen candidates.\lrfootnote{See \casecite[1--2]{mobleyHiredScore2025}.} The order managed collective notice rather than applying this Article's statute, but it illustrates a similar operational sequence: identify deployment and output, remove customers with no relevant run or use, and reserve disputed effect.

Stage One generalizes that logic at the event level. The employer may control job criteria, feature activation, decision weight, review, notice, and disposition; the vendor may control design, defaults, validation, updates, and cross-customer monitoring. A customer exits the control lane if no relevant feature ran, no output was produced, or no output reached the decision path. Generic hosting, a post-event acquisition, or transmission of a fixed customer rule can support the same result. A surviving lane then defines Stage Two: aggregate recommendation and validation data for disparate impact; decision inputs and review for disparate treatment; or access, exceptions, and review for accommodation, each bounded by the governing law and procedure.

Where hosting and legal control remain divided, a coordinated protocol replaces serial referrals. The employer authorizes the defined customer fields; the platform runs the query against the correct tenant and version; both verify the schema; direct identifiers are pseudonymized; any lawful protected-characteristic linkage remains segregated; and experts use a secure environment. The protocol identifies which enterprise supplies each field and who attests to completeness. It neither transfers all customer data nor treats physical hosting as universal Rule 34 control.

\subsubsection{A Stable Public-Employment Analogue}

\emph{Houston Federation of Teachers v. Houston Independent School District} supplies a completed public-employment analogue. HISD licensed SAS's proprietary teacher-evaluation system, acknowledged that it did not verify the calculations and that teachers could not reproduce them, and had treated requested materials as outside its custody or control.\lrfootnote{See \casecite[1172--73, 1177--78]{houstonTeachers2017}.} The court denied HISD summary judgment on procedural due process without imposing liability on SAS or ordering public release of source code; it explained that relief could restrict government use of an unverifiable procedure while preserving the trade secret.\lrfootnote{See \casecite[1179--83, 1185]{houstonTeachers2017}.} The front door supplies the prospective complement: define version, input, validation, audit, and challenge records before adoption, then allocate protected access without conflating public duty, vendor secrecy, or ultimate liability.
